\scanpage{203}
En plus de la \uline{structure entière} sur le \deleted{\unsure{él}} \inserted{foncteur} de Hodge-DeRham, relatif à un pt complexe, il faut signaler la \uline{décomposition en types}, également de nature transcendante (elle est compatible avec la filtration canonique). On peut aussi dire que cette dernière \deleted{\unsure{associée}} \inserted{est déterminée via filtration} par \uline{l’opération involutive} \inserted{$\tau$ $\tfrac12$ linéaire} $\tau'$ sur $\unsure{\xi|\zeta}_{x,\infty}$, provenant de l’écriture

\[
\xi_{x,\infty}=\xi_{x,\mathbb{R}}\otimes_{\mathbb{R}}\mathbb{C}
\qquad\text{où}\quad \xi_{x,\mathbb{R}}\simeq\xi_{x,\mathbb{Z}}\otimes_{\mathbb{Z}}\mathbb{R}.
\]

Si $k(a)$ est réel, on a une autre application $\tfrac12$ linéaire involutive $\tau''$, provenant de

\[
\xi_{\unsure{\alpha|x},\infty}\simeq\xi_{a,\infty}{}_{\mathbb{R}}\otimes_{\mathbb{R}}\mathbb{C}
\]

et $\tau$ sur $(\xi_{x,\mathbb{Z}})_{\mathbb{C}}$ n’est autre que $\tau'\tau''$

\diagram{\begin{tikzcd}[ampersand replacement=\&,column sep=0.3em,row sep=1.2ex]
{} \& \tau={} \& \tau' \arrow[dl,dash,bend right=25] \& \tau'' \arrow[dr,dash,bend right=20] \& {}=\tau''\tau' \\
{} \& \substack{\tfrac12\ \text{linéaire}\\\text{transcendant}} \& {} \& {} \& \substack{\tfrac12\ \text{linéaire}\\\text{algébrique}}
\end{tikzcd}}

% slanted
\margin{à vérifier}

La situation locale \inserted{(dans le cas réel)} est connue par la connaissance de $\xi_{\unsure{x|a},\infty}^{\tau''}=(\xi_{a,\infty})_{\mathbb{R}}$ \inserted{(Hodge-DeRham réel)}, muni de $\tau=\tau'$ (involution transcendante, qui est $\mathbb{R}$-linéaire sur le dit espace), ou \inserted{mieux par celle de $\tau=\tau''$ opérant sur $\xi_{x,\infty}^{\tau'}\simeq\xi_{x,\mathbb{R}}$,} \uline{\unsure{évidemt} de $\tau$ sur $\xi_{a,\mathbb{Z}}$}, \unsure{donc au choix de|d’où on déduit de|dont on déduit de}

\[
\boxed{\tau\ \text{sur}\ \xi_{a,\ell}}
\]

[or c’est là \inserted{l’opération de} l’élément de $\unsure{\Pi|\pi}_1(X)$, \uline{frobénius} en la place \deleted{\unsure{f}} réelle - !]

% slanted
\margin{à vérifier}

% archive sheet number: 202
